\textbf{概率与期望速查}

\textbf{适用条件}
除特别说明外，随机变量可积；涉及方差时要求二阶矩有限。
$[P]$ 为事件 $P$ 的指示变量。独立性只在明确写出时使用。

\textbf{指示变量与线性性}
\[
\operatorname{E}[X]=\sum_i\operatorname{E}[X_i]\quad(X=\sum_iX_i),\qquad
\operatorname{E}[[P]]=\Pr(P).
\]
求“满足某性质的对象个数”时，给每个候选对象建指示变量再求和；
各事件无须独立。
例如均匀随机排列 $n\ge1$ 个元素，固定点期望为
$n\cdot(1/n)=1$；逆序对期望为 $\binom n2/2$。

\textbf{整数等待时间的尾和公式}
若 $T$ 为非负整数值随机变量，
\[
\operatorname{E}[T]=\sum_{k\ge1}\Pr(T\ge k)
=\sum_{k\ge0}\Pr(T>k),
\]
右侧允许取 $+\infty$。几何等待时间 $T\in\{1,2,\ldots\}$，
每轮独立成功概率 $p\in(0,1]$ 时，
$\Pr(T>k)=(1-p)^k$，故 $\operatorname{E}[T]=1/p$。

\textbf{条件期望与方差}
若 $A_1,\ldots,A_m$ 是概率为正的有限划分，则
\[
\operatorname{E}[X]=\sum_{i=1}^m\Pr(A_i)\operatorname{E}[X\mid A_i].
\]
\[
\begin{aligned}
\operatorname{Var}(X)&=\operatorname{E}[X^2]-\operatorname{E}[X]^2,\\
\operatorname{Var}\!\left(\sum_iX_i\right)
&=\sum_i\operatorname{Var}(X_i)
+2\sum_{i<j}\operatorname{Cov}(X_i,X_j).
\end{aligned}
\]
相互独立或两两协方差为 $0$ 时，方差可直接相加；一般需保留协方差项。

\textbf{集齐所有种类的期望时间}
每次独立抽取 $n\ge1$ 种之一，均匀分布时，
\[
\operatorname{E}[T]=nH_n=n\sum_{k=1}^{n}\frac1k.
\]
若第 $i$ 种概率 $p_i>0$，$\sum_i p_i=1$，则容斥给出
\[
\operatorname{E}[T]=
\sum_{\emptyset\ne S\subseteq[n]}
(-1)^{|S|+1}\frac1{\sum_{i\in S}p_i}.
\]
该式枚举 $2^n-1$ 个非空子集；概率为 $0$ 的必需种类会使期望无穷大。
